On considère l’équation différentielle (E) : $y' + 4y = 8-3e^{-4x}$ où $y$ est une fonction de la variable réelle $x$, définie et dérivable sur l’ensemble $\R$ des nombres réels, et $y'$ est la fonction dérivée de $y$.

\begin{enumerate}
	\item Résoure l'équation homogène (E$_0$) : $y'+4y=0$\sautp

	\Lignes{2}
	\item A l’aide du logiciel Géogébra, déterminer les nombres réels $a$ et $b$ pour que la fonction $p$ définie sur $\R$ par : $p(x) = a + b~x~e^{-4x}$ soit une solution particulière de l'équation différentielle (E), et donner l’expression de $p$. \sautm
	
	On trouve : $a=$\makebox[0.1\linewidth]{\dotfill} \hfill $b=$\makebox[0.1\linewidth]{\dotfill} \hfill $p(x)=$\makebox[0.3\linewidth]{\dotfill}\sauts
	\item En déduire la forme générale des solutions de (E).\sautp

	\Lignes{1}
	\item Déterminer, à l'aide de géogebra et d'un curseur bien choisi, la solution particulière $f$ de (E) telle que $f(0)=2$.\sautp

	\Lignes{1}
\end{enumerate}